\documentclass[11pt,a4paper]{article}

\usepackage[T1]{fontenc}
\usepackage{mlmodern}
\usepackage[margin=2.2cm]{geometry}
\usepackage{amsmath,amssymb}
\usepackage{booktabs}
\usepackage{enumitem}
\usepackage{xcolor}
\usepackage{listings}
\usepackage{algorithm}
\usepackage{algpseudocode}
\usepackage[colorlinks=true,linkcolor=black,urlcolor=black,citecolor=black,filecolor=black]{hyperref}

\definecolor{codebg}{RGB}{247,247,249}
\definecolor{codekw}{RGB}{0,70,140}
\definecolor{codecm}{RGB}{110,110,110}
\definecolor{codestr}{RGB}{160,40,40}

\lstset{
  language=Python,
  basicstyle=\ttfamily\footnotesize,
  keywordstyle=\color{codekw}\bfseries,
  commentstyle=\color{codecm}\itshape,
  stringstyle=\color{codestr},
  backgroundcolor=\color{codebg},
  frame=single, rulecolor=\color{black!15},
  breaklines=true, columns=fullflexible, keepspaces=true,
  showstringspaces=false, xleftmargin=4pt, xrightmargin=4pt,
  aboveskip=6pt, belowskip=6pt
}

\setlist{itemsep=2pt,topsep=3pt}
\newcommand{\E}{\mathbb{E}}
\newcommand{\code}[1]{\texttt{#1}}

\title{\textbf{Adaptive Curriculum and Self-Play for Generalizable Adversarial RL Agents}\\[4pt]
\large Theory and Implementation Study Guide}
\author{Aayush Shukla \and Aditya Kumar}
\date{}

\begin{document}
\maketitle
\tableofcontents
\newpage

%=====================================================================
\section{The Problem in One Paragraph}
A \textbf{Hunter} and an \textbf{Evader} play on procedurally generated $21\times21$ dungeons.
The Evader must find a key, open a locked door and reach the exit; the Hunter must catch it.
Both move simultaneously (UP, DOWN, LEFT, RIGHT, STAY), see only a $9\times9$ line-of-sight
window, and an episode lasts at most 200 steps. We train both agents with reinforcement learning
(PPO and DQN, implemented from scratch on PyTorch tensors and autograd, no RL libraries) and ask:
\emph{does structured training (curriculum + self-play) produce agents that generalize to unseen
maps and are harder to exploit than naive simultaneous training?}

\paragraph{Observation (496 numbers).} A $9\times9$ window with 6 binary channels
(wall, door, key, exit, other agent, visible) $= 486$ values, plus 10 scalars (own position,
key held, time fraction, direction and distance to target, other-agent visibility and
relative position).

\paragraph{Reward.} Win $+1$, loss $-1$, step penalty $-0.01$, shaping $0.01$ per tile of
path-distance progress (Evader toward its objective, Hunter toward the Evader), Evader $+0.1$
for the key and $+0.1$ for opening a door. The win/loss part is zero-sum; the shaping is not.

%=====================================================================
\section{Reinforcement Learning Foundations}

\subsection{MDP, POMDP and Markov game}
\begin{itemize}
  \item \textbf{MDP} $(\mathcal{S},\mathcal{A},P,R,\gamma)$: state $s_t$, action $a_t\sim\pi(\cdot|s_t)$,
        next state $s_{t+1}\sim P(\cdot|s_t,a_t)$, reward $r_t$.
  \item \textbf{POMDP}: the agent sees an observation $o_t$ (our $9\times9$ window), not the full state.
        A memoryless policy $\pi(a|o)$ may then need to be \emph{stochastic} to be optimal.
  \item \textbf{Markov game} (stochastic game): several agents act at once; each agent's transition
        and reward depend on the joint action. Ours is a two-player, (mostly) zero-sum, partially
        observable Markov game.
\end{itemize}

\subsection{Return, value functions, advantage}
\[
G_t=\sum_{k=0}^{\infty}\gamma^k r_{t+k},\qquad
V^\pi(s)=\E_\pi[G_t\mid s_t=s],\qquad
Q^\pi(s,a)=\E_\pi[G_t\mid s_t=s,a_t=a],
\]
\[
A^\pi(s,a)=Q^\pi(s,a)-V^\pi(s)\quad\text{(how much better $a$ is than the policy's average).}
\]
Bellman equations:
$V^\pi(s)=\E_{a\sim\pi}\big[r+\gamma V^\pi(s')\big]$,\;
$Q^*(s,a)=\E\big[r+\gamma\max_{a'}Q^*(s',a')\big]$.
We use $\gamma=0.99$ (effective horizon $\approx 1/(1-\gamma)=100$ steps).

\subsection{Two families of methods}
\begin{center}
\begin{tabular}{lll}
\toprule
 & \textbf{Value-based (DQN)} & \textbf{Policy-gradient (PPO)}\\
\midrule
Learns & $Q(s,a)$, acts $\arg\max_a Q$ & $\pi_\theta(a|s)$ directly (+ critic $V$)\\
Data & off-policy (replay buffer) & on-policy (fresh rollouts)\\
Policy & deterministic ($\varepsilon$-greedy to explore) & stochastic by construction\\
Stability & target network, replay & clipped updates\\
Sample efficiency & higher (reuses data) & lower, but simpler and more stable\\
\bottomrule
\end{tabular}
\end{center}

\subsection{Policy gradient theorem}
For $J(\theta)=\E_{\pi_\theta}[G_0]$:
\[
\nabla_\theta J(\theta)=\E_{\pi_\theta}\big[\nabla_\theta\log\pi_\theta(a_t|s_t)\,A^{\pi}(s_t,a_t)\big].
\]
Intuition: increase the log-probability of actions that turned out better than average,
decrease it for worse ones. REINFORCE uses $G_t$ in place of $A$ (high variance); subtracting a
baseline $V(s)$ keeps the gradient unbiased but reduces variance $\Rightarrow$ \textbf{actor-critic}.

%=====================================================================
\section{Generalized Advantage Estimation (GAE)}
The one-step TD error is $\delta_t=r_t+\gamma V(s_{t+1})(1-d_t)-V(s_t)$, where $d_t=1$ if the
episode ended after step $t$. GAE averages $n$-step estimates with weight $\lambda$:
\[
\hat A_t=\sum_{l=0}^{\infty}(\gamma\lambda)^l\delta_{t+l}
\quad\Longleftrightarrow\quad
\hat A_t=\delta_t+\gamma\lambda(1-d_t)\hat A_{t+1}.
\]
$\lambda=0$: one-step TD (low variance, biased). $\lambda=1$: Monte-Carlo (unbiased, high variance).
We use $\lambda=0.95$. Value targets are $\hat R_t=\hat A_t+V(s_t)$.

\paragraph{Timeouts.} Treating a 200-step timeout as terminal is normally a bug (the state is
not really terminal), but our observation contains the time fraction $t/200$, so time is part of the
state and the episode end \emph{is} Markov. That is why we do not bootstrap on truncation.

\begin{lstlisting}[title={\code{rl/ppo.py} -- RolloutBuffer.compute\_gae (exact)}]
def compute_gae(self, last_values, gamma, lam):
    # dones[t] = the episode ended after step t, so step t does not bootstrap.
    # Timeouts are treated as terminal: the time fraction is part of the observation.
    adv = np.zeros_like(self.rewards)
    gae = np.zeros(self.N, np.float32)
    next_v = last_values
    for t in reversed(range(self.T)):
        nonterm = 1.0 - self.dones[t]
        delta = self.rewards[t] + gamma * next_v * nonterm - self.values[t]
        gae = delta + gamma * lam * nonterm * gae
        adv[t] = gae
        next_v = self.values[t]
    return adv, adv + self.values
\end{lstlisting}
The buffer has shape $(T,N)$: $T=256$ steps for each of $N=32$ parallel envs, i.e.\ 8192 samples per update.

%=====================================================================
\section{Proximal Policy Optimization (PPO)}

\subsection{Why PPO exists}
Vanilla policy gradient is very sensitive to step size: one large update can collapse the policy,
and because the data is on-policy a bad policy then collects bad data. \textbf{TRPO} fixes this by
constraining $\mathrm{KL}(\pi_{\text{old}}\,\|\,\pi_\theta)\le\delta$, but needs second-order
optimisation. \textbf{PPO} (Schulman et al., 2017) gets a similar ``trust region'' effect with a
simple clipped first-order objective, so it can safely do \emph{several epochs} of minibatch SGD on
the same rollout.

\subsection{The clipped surrogate objective}
Probability ratio $\;r_t(\theta)=\dfrac{\pi_\theta(a_t|s_t)}{\pi_{\theta_\text{old}}(a_t|s_t)}=\exp(\log\pi_\theta-\log\pi_{\text{old}})$.
\[
L^{\text{CLIP}}(\theta)=\E_t\Big[\min\big(r_t\hat A_t,\;\operatorname{clip}(r_t,1-\epsilon,1+\epsilon)\,\hat A_t\big)\Big],\qquad \epsilon=0.2.
\]
\begin{itemize}
  \item If $\hat A_t>0$ (good action): the objective grows with $r_t$ only up to $1+\epsilon$; beyond that
        the gradient is zero, so we cannot push the probability up too far in one update.
  \item If $\hat A_t<0$ (bad action): it is capped at $1-\epsilon$, so we cannot push it down too far.
  \item The $\min$ makes it a \emph{pessimistic} lower bound: clipping only ever removes incentive,
        never adds it.
\end{itemize}

\subsection{Full loss}
\[
\mathcal{L}(\theta)=-L^{\text{CLIP}}
+c_v\,\tfrac12\,\E_t\big[(V_\theta(s_t)-\hat R_t)^2\big]
-c_e\,\E_t\big[\mathcal{H}(\pi_\theta(\cdot|s_t))\big],
\qquad c_v=0.5,\;c_e=0.01.
\]
The entropy bonus $\mathcal{H}=-\sum_a\pi\log\pi$ keeps the policy exploring.

\subsection{Implementation details we use (and why)}
\begin{center}
\begin{tabular}{lll}
\toprule
\textbf{Detail} & \textbf{Value} & \textbf{Why}\\
\midrule
Networks & separate actor and critic MLPs, 256-256, tanh & no interference between losses\\
Init & orthogonal; gain $\sqrt2$, actor head $0.01$, critic $1$ & near-uniform initial policy\\
Rollout & 32 envs $\times$ 256 steps $=8192$ samples & lower-variance gradient\\
Epochs / minibatch & 4 epochs, minibatch 512 & data reuse\\
Advantage norm. & per minibatch, $(\hat A-\mu)/(\sigma+10^{-8})$ & scale-free updates\\
KL early stop & stop if approx-KL $>1.5\times0.03$ & extra trust-region safety\\
LR & $3\cdot10^{-4}$, Adam ($\epsilon{=}10^{-5}$), linear anneal to 10\% & stable late training\\
Grad clipping & global norm $0.5$ & avoids exploding updates\\
\bottomrule
\end{tabular}
\end{center}
\textbf{Approx.\ KL} uses the estimator $\E[(r-1)-\log r]$, which is always $\ge0$ and low-variance.
\textbf{Clip fraction} is the share of samples with $|r-1|>\epsilon$.
\textbf{Explained variance} $1-\mathrm{Var}(\hat R-V)/\mathrm{Var}(\hat R)$ measures critic quality
($1$ = perfect, $\le0$ = useless).

\subsection{PPO algorithm}
\begin{algorithm}[h]
\caption{PPO (as implemented)}
\begin{algorithmic}[1]
\For{update $=1,\dots,U$}
  \State Run $\pi_{\theta}$ in $N$ envs for $T$ steps; store $o_t,a_t,\log\pi_\text{old}(a_t|o_t),V(o_t),r_t,d_t$
  \State Compute $\hat A_t,\hat R_t$ with GAE using $V(o_{T})$ as bootstrap
  \State Set learning rate $\alpha\leftarrow\alpha_0\max(0.1,\,1-\text{progress})$
  \For{epoch $=1,\dots,K$}
    \For{each shuffled minibatch $B$}
      \State $r\leftarrow\exp(\log\pi_\theta(a|o)-\log\pi_\text{old})$, normalise $\hat A$ over $B$
      \State $\mathcal{L}\leftarrow -L^\text{CLIP}+c_v L^V-c_e\mathcal{H}$; Adam step with grad-norm clip
      \If{approx-KL $>1.5\,\text{KL}_\text{target}$} \textbf{stop all epochs} \EndIf
    \EndFor
  \EndFor
\EndFor
\end{algorithmic}
\end{algorithm}

\subsection{Code: network and action sampling}
\begin{lstlisting}[title={\code{rl/nets.py} -- mlp, ActorCritic, sample\_categorical (exact)}]
def mlp(sizes, act=nn.Tanh, out_gain=1.0):
    layers = []
    for i in range(len(sizes) - 1):
        lin = nn.Linear(sizes[i], sizes[i + 1])
        last = i == len(sizes) - 2
        nn.init.orthogonal_(lin.weight, out_gain if last else np.sqrt(2))
        nn.init.zeros_(lin.bias)
        layers.append(lin)
        if not last:
            layers.append(act())
    return nn.Sequential(*layers)


class ActorCritic(nn.Module):
    """Separate policy and value MLPs (no shared trunk)."""

    def __init__(self, obs_dim, n_actions, hidden=(256, 256)):
        super().__init__()
        self.obs_dim, self.n_actions, self.hidden = obs_dim, n_actions, list(hidden)
        self.pi = mlp([obs_dim, *hidden, n_actions], out_gain=0.01)
        self.v = mlp([obs_dim, *hidden, 1], out_gain=1.0)

    def logits(self, obs):
        return self.pi(obs)

    def value(self, obs):
        return self.v(obs).squeeze(-1)


def sample_categorical(logits, deterministic=False, generator=None):
    """Returns actions and their log-probabilities."""
    logp_all = torch.log_softmax(logits, dim=-1)
    if deterministic:
        a = logits.argmax(-1)
    else:
        a = torch.multinomial(logp_all.exp(), 1, generator=generator).squeeze(-1)
    return a, logp_all.gather(-1, a.unsqueeze(-1)).squeeze(-1)
\end{lstlisting}

\subsection{Code: the PPO update (core loop, exact)}
\begin{lstlisting}[title={\code{rl/ppo.py} -- PPO.update (statistics bookkeeping omitted)}]
def update(self, buf: RolloutBuffer, last_values, progress=0.0):
    c = self.cfg
    lr = c.lr * max(c.min_lr_frac, 1.0 - progress) if c.anneal_lr else c.lr
    for g in self.opt.param_groups:
        g["lr"] = lr

    adv, ret = buf.compute_gae(last_values, c.gamma, c.gae_lambda)
    n = buf.T * buf.N
    obs = self._t(buf.obs.reshape(n, -1))
    act = self._t(buf.actions.reshape(n))
    old_logp = self._t(buf.logp.reshape(n))
    adv_t = self._t(adv.reshape(n))
    ret_t = self._t(ret.reshape(n))

    for _ in range(c.epochs):
        perm = torch.randperm(n, device=self.device)
        for s in range(0, n, c.minibatch):
            idx = perm[s:s + c.minibatch]
            logp_all = torch.log_softmax(self.net.logits(obs[idx]), dim=-1)
            logp = logp_all.gather(-1, act[idx].unsqueeze(-1)).squeeze(-1)
            log_ratio = logp - old_logp[idx]
            ratio = log_ratio.exp()

            a = adv_t[idx]
            a = (a - a.mean()) / (a.std() + 1e-8)
            pg_loss = -torch.min(ratio * a,
                                 ratio.clamp(1 - c.clip, 1 + c.clip) * a).mean()
            v_loss = 0.5 * (self.net.value(obs[idx]) - ret_t[idx]).pow(2).mean()
            entropy = -(logp_all.exp() * logp_all).sum(-1).mean()
            loss = pg_loss + c.vf_coef * v_loss - c.ent_coef * entropy

            self.opt.zero_grad()
            loss.backward()
            torch.nn.utils.clip_grad_norm_(self.net.parameters(), c.max_grad_norm)
            self.opt.step()

            with torch.no_grad():
                kl = ((ratio - 1) - log_ratio).mean().item()
            if c.target_kl is not None and kl > 1.5 * c.target_kl:
                ...  # early stop: break out of both loops
\end{lstlisting}

\subsection{Code: collecting a rollout}
\begin{lstlisting}[title={\code{rl/ppo.py} -- PPO.act (exact)}]
@torch.no_grad()
def act(self, obs, deterministic=False):
    o = self._t(obs)
    a, logp = sample_categorical(self.net.logits(o), deterministic)
    v = self.net.value(o)
    return a.cpu().numpy(), logp.cpu().numpy(), v.cpu().numpy()
\end{lstlisting}
\begin{lstlisting}[title={\code{training/train.py} -- training loop skeleton (abridged)}]
for u in range(n_updates):
    for _ in range(T):                                   # T = 256 steps
        actions = np.full((N, 2), -1, np.int64)          # -1 = scripted inside worker
        for r in roles:                                  # learners act
            idx = learner_envs[r]
            a, logp, v = learners[r].act(obs[idx, ROLES.index(r)])
            actions[idx, ROLES.index(r)] = a
        ...                                              # frozen opponents act
        obs, rew, done, eps = venv.step(actions)         # 8 worker processes
        for r in roles:
            bufs[r].add(o, a, logp, v, rew[idx, ROLES.index(r)], done[idx])
    for r in roles:
        last_v = learners[r].value(obs[idx, ROLES.index(r)])
        learners[r].update(bufs[r], last_v, progress)
    # curriculum.update(...), pool snapshot (self-play), periodic eval + checkpoint
\end{lstlisting}

%=====================================================================
\section{Deep Q-Network (DQN)}

\subsection{From Q-learning to DQN}
Tabular Q-learning: $Q(s,a)\leftarrow Q(s,a)+\alpha\big[r+\gamma\max_{a'}Q(s',a')-Q(s,a)\big]$.
DQN (Mnih et al., 2015) replaces the table with a neural net $Q_\phi$ and minimises
$\big(Q_\phi(s,a)-y\big)^2$. Naively this diverges (``deadly triad'': function approximation +
bootstrapping + off-policy data). DQN stabilises it with two tricks:
\begin{enumerate}
  \item \textbf{Experience replay}: store transitions in a buffer (ours: 200k, stored as float16) and
        train on random minibatches, which breaks temporal correlation and reuses data.
  \item \textbf{Target network} $Q_{\bar\phi}$: a frozen copy used for targets, synced every
        2000 gradient steps, so the target does not move with every update.
\end{enumerate}

\subsection{Double DQN, Huber loss, exploration}
Plain DQN target $y=r+\gamma(1-d)\max_{a'}Q_{\bar\phi}(s',a')$ \emph{overestimates} values
(max of noisy estimates is biased upward). \textbf{Double DQN} decouples selection and evaluation:
\[
y=r+\gamma(1-d)\,Q_{\bar\phi}\big(s',\arg\max_{a'}Q_\phi(s',a')\big).
\]
The \textbf{Huber} (smooth-L1) loss is quadratic for small errors and linear for large ones,
so outliers do not produce huge gradients. Exploration is $\varepsilon$-greedy with $\varepsilon$
decaying linearly from $1.0$ to $0.05$ over the first 30\% of training. Learning starts after
20k samples; lr $2.5\cdot10^{-4}$, batch 256.

\begin{lstlisting}[title={\code{rl/dqn.py} -- DQN.act and DQN.train\_step (exact)}]
@torch.no_grad()
def act(self, obs, eps):
    a = self.q(torch.as_tensor(obs, device=self.device)).argmax(-1).cpu().numpy()
    rand = self.rng.random(len(a)) < eps
    a[rand] = self.rng.integers(self.n_actions, size=int(rand.sum()))
    return a

def train_step(self):
    c = self.cfg
    o, a, r, o2, d = (torch.as_tensor(x, device=self.device)
                      for x in self.buf.sample(c.batch_size))
    q = self.q(o).gather(-1, a.unsqueeze(-1)).squeeze(-1)
    with torch.no_grad():
        if c.double:
            a2 = self.q(o2).argmax(-1, keepdim=True)
            q2 = self.target(o2).gather(-1, a2).squeeze(-1)
        else:
            q2 = self.target(o2).max(-1).values
        y = r + c.gamma * (1 - d) * q2
    loss = F.smooth_l1_loss(q, y)
    self.opt.zero_grad()
    loss.backward()
    torch.nn.utils.clip_grad_norm_(self.q.parameters(), c.max_grad_norm)
    self.opt.step()
    self.n_grad += 1
    if self.n_grad % c.target_every == 0:
        self.target.load_state_dict(self.q.state_dict())
    return loss.item(), q.mean().item()
\end{lstlisting}

\paragraph{Why PPO is our main algorithm.} (1) In multi-agent training the opponent changes, so
old replay data describes an outdated opponent; on-policy PPO always learns from current games.
(2) Under partial observability, stochastic policies can be strictly better than any deterministic
one; PPO learns a stochastic policy natively, DQN's $\arg\max$ is deterministic.
(3) PPO is robust to hyper-parameters. DQN is our value-based baseline and gives real $Q$-values
for the heatmaps.

%=====================================================================
\section{Multi-Agent Training: Independent Learning vs Self-Play}

\subsection{Non-stationarity}
From one agent's viewpoint, the opponent is part of the environment. If the opponent is also
learning, the environment's dynamics change every update, so the Markov assumption behind
single-agent RL is broken. Consequences: oscillation, \textbf{cycling} (A beats B, B adapts,
A forgets how to beat older B) and \textbf{catastrophic forgetting}.

\subsection{Our three training modes}
\begin{description}
  \item[single] one PPO learner vs a fixed opponent (scripted, or a frozen trained model).
        Stationary, so it is the baseline.
  \item[independent] both roles learn \emph{simultaneously} from the same games, each treating
        the other as environment (``independent learners''). Naive, non-stationary.
  \item[selfplay] each role learns against a \textbf{pool of frozen snapshots} of the other role.
        Half of the 32 envs train the Hunter, half train the Evader.
\end{description}
\texttt{--steps} counts samples \emph{per learner}, so every mode gives each learner the same data.

\subsection{Self-play with an opponent pool}
\begin{itemize}
  \item Every 50k samples, a frozen copy of each learner is added to its pool (max 20; oldest dropped).
  \item At the start of each episode an env samples its opponent: the latest snapshot with
        probability $0.5$, otherwise uniformly from the pool.
  \item Playing the latest keeps pressure on current weaknesses; playing older ones prevents
        forgetting and cycling. This is related to \emph{fictitious play}, where one best-responds to
        the \emph{average} of past opponent strategies, which converges to a Nash equilibrium in
        two-player zero-sum games. AlphaStar's league and OpenAI Five use the same idea.
  \item The opponent is frozen during an episode, so each learner faces a stationary distribution
        of opponents between snapshots.
\end{itemize}
\begin{lstlisting}[title={\code{rl/pool.py} -- OpponentPool (exact)}]
class OpponentPool:
    """Frozen snapshots of one role. Sampling: the latest snapshot with probability
    `p_latest`, otherwise uniform over the whole pool. Oldest snapshots are dropped."""

    def __init__(self, max_size=20, p_latest=0.5, seed=0):
        self.items, self.max_size, self.p_latest = [], max_size, p_latest
        self.rng = np.random.default_rng(seed)

    def add(self, net, tag):
        self.items.append(FrozenPolicy(net, tag))
        if len(self.items) > self.max_size:
            self.items.pop(0)

    def sample(self):
        if self.rng.random() < self.p_latest:
            return self.items[-1]
        return self.items[int(self.rng.integers(len(self.items)))]
\end{lstlisting}

\subsection{Measuring stability: cross-play}
We save checkpoints of both roles during training and play every Hunter checkpoint against every
Evader checkpoint. Entry $(i,j)$ is the Hunter win rate. In stable training a later Hunter beats
all earlier Evaders, giving a monotone matrix; cycling shows up as non-transitive bands. We
report a \emph{regression rate}: the fraction of (later, earlier) Hunter pairs where the later one
is worse by more than $0.1$ against some Evader. Other stability metrics: the standard deviation of
the validation win-rate curve over the second half of training, and the largest drop between
consecutive evaluations.

%=====================================================================
\section{Curriculum Learning}
\textbf{Idea} (Bengio et al., 2009): learn easy tasks first, then harder ones, the way humans
are taught. In RL, the main benefit is against \textbf{sparse rewards}. On a full map a random
agent almost never wins, so there is little signal. On a small map it wins often enough to learn
basic skills, which transfer.

\subsection{Our levels}
Map complexity and scripted-opponent skill increase together
(\code{opp\_eps} = probability the scripted opponent takes a random action).
\begin{center}
\begin{tabular}{cccccc}
\toprule
Level & Size & Rooms & Doors & Obstacles & \code{opp\_eps}\\
\midrule
0 & $13\times13$ & 3 & 0 & 0.0 & 0.75\\
1 & $15\times15$ & 3 & 1 & 0.1 & 0.50\\
2 & $17\times17$ & 4 & 1 & 0.1 & 0.25\\
3 (target) & $21\times21$ & 5 & 1 & 0.2 & 0.00\\
\bottomrule
\end{tabular}
\end{center}
The observation is a fixed $9\times9$ egocentric window, so the same network works on every map
size, which is what makes a size curriculum possible.

\subsection{Adaptive rule}
\begin{itemize}
  \item \textbf{Single learner}: advance when its win rate over the last 300 episodes is $\ge0.6$.
  \item \textbf{Two learners}: win rate is relative (it hovers around a balance), so instead we use
        the \emph{decisive rate}: the share of games that end in a catch or an escape rather than a
        timeout, $\ge0.8$. Timeouts mean neither agent is competent yet.
  \item At least 50k samples per level; fallback: advance anyway after 20\% of the budget.
  \item Modes \code{linear} (fixed schedule over the first 60\%) and \code{off} are the comparisons.
\end{itemize}
\begin{lstlisting}[title={\code{rl/curriculum.py} -- Curriculum.update (exact)}]
def update(self, timesteps):
    """Returns True if the level changed."""
    if self.final or self.mode == "off":
        return False
    if self.mode == "linear":
        per = self.ramp * self.total / (len(self.levels) - 1)
        target = min(len(self.levels) - 1, int(timesteps // per))
        advance = target > self.level
    else:
        at_level = timesteps - self.level_start
        full = len(self.recent) == self.recent.maxlen
        advance = (at_level >= self.min_steps and full and self.rate() >= self.threshold) \
            or at_level >= self.max_steps
    if advance:
        self.level += 1
        self.level_start = timesteps
        self.recent.clear()
        self.history.append((timesteps, self.level))
    return advance
\end{lstlisting}

%=====================================================================
\section{Generalization and Procedural Generation}
Evaluating on training maps measures \emph{memorisation}. Procedural content generation (PCG)
lets us create an unlimited number of distinct maps from integer seeds (Cobbe et al., Procgen, 2019),
so generalization can be measured directly:
\begin{center}
\begin{tabular}{lll}
\toprule
Split & Seeds & Use\\
\midrule
train & 0--999 & training only\\
val & 10000--10199 & periodic evaluation during training\\
test & 20000--20499 & final numbers only, never tuned on\\
OOD & test seeds, $25\times25$, 7 rooms, 2 doors & out-of-distribution stress test\\
\bottomrule
\end{tabular}
\end{center}
\textbf{Generalization gap} $=$ win(train) $-$ win(test). A large positive gap means overfitting to
layouts. The OOD drop $=$ win(test) $-$ win(OOD) measures robustness to larger, more complex maps.
The generator (rooms + L-shaped corridors + a door on a cut corridor + key in the reachable region)
validates that every map is solvable.

%=====================================================================
\section{Exploitability}
In game theory, the exploitability of a strategy $\pi$ is how much a \emph{best response} to it
gains: $\text{expl}(\pi)=\max_{\pi'}J(\pi',\pi)-J^\star$. A Nash-equilibrium strategy has
exploitability 0. Exact best responses are intractable here, so we \textbf{approximate the best
response by training a fresh PPO agent} against the frozen model and measure:
\begin{itemize}
  \item \textbf{samples to threshold}: samples until the fresh learner's rolling win rate (200 episodes)
        reaches 50\%. Fewer samples $\Rightarrow$ more exploitable;
  \item the final win rate and area under the learning curve against the frozen model.
\end{itemize}
Hypothesis: agents from self-play (which have faced many opponents) are harder to exploit than
agents trained against one fixed opponent or with naive independent learning.

%=====================================================================
\section{Explainability: Decision Heatmaps}
Win rate says \emph{whether} an agent works, not \emph{why}. For a fixed map and opponent position
we place the agent on every walkable cell, compute its observation, and plot:
\begin{itemize}
  \item \textbf{Value map} $V(s)$ (PPO critic) or $\max_aQ(s,a)$ (DQN): where the agent believes it is
        winning. A good Hunter's value rises near the exit it guards; a good Evader's value rises
        along the path to key, door and exit.
  \item \textbf{Policy arrows}: the most likely action per cell, length scaled by its probability.
  \item \textbf{Entropy map} $\mathcal{H}(\pi(\cdot|s))$: where the agent is uncertain. High entropy
        appears where the opponent is out of sight.
\end{itemize}
The live viewer also shows action probabilities as arrows around each agent in real time.

%=====================================================================
\section{Key Empirical Finding: Stochastic vs Greedy Policies}
Evaluating the PPO Hunter (1M samples vs the scripted path Evader, 200 maps):
\begin{center}
\begin{tabular}{lccc}
\toprule
Policy at test time & val win & test win & timeouts\\
\midrule
greedy ($\arg\max$) & 0.475 & 0.45 & 30--39\%\\
sampled from $\pi$ & 0.68 & 0.705 & 9--12\%\\
\bottomrule
\end{tabular}
\end{center}
\textbf{Explanation.} The policy is memoryless and the world is partially observable. When the
Hunter loses sight of the Evader, many different states look identical, and a deterministic policy
does the same thing in all of them. It oscillates (e.g.\ left--right) and runs out the clock.
Sampling breaks these loops. This matches the theory: in POMDPs the best memoryless policy can be
stochastic (Singh, Jaakkola \& Jordan, 1994). We therefore evaluate PPO by sampling, with a seeded
generator so results are reproducible. A recurrent (LSTM) policy with memory would be the
principled fix.

%=====================================================================
\section{Engineering Notes}
\begin{itemize}
  \item \textbf{Vectorized envs}: 32 envs in 8 worker processes (\code{multiprocessing} pipes);
        scripted opponents run inside the workers, and policies run batched in the main process.
        Throughput is $\sim$5.5--8k env steps/s, versus $\sim$475 with the earlier library-based version.
  \item \textbf{Why a GPU barely helps}: about 75\% of time is the Python environment simulation
        (line-of-sight per cell, BFS distances), which runs only on the CPU. The network is tiny
        ($\approx$400k parameters), and the rollout loop is sequential (policy $\to$ env $\to$ policy).
        By Amdahl's law, speeding up the remaining $\sim$25\% (updates) infinitely gives at most
        $1/0.75\approx1.33\times$; realistically $1.1$--$1.2\times$.
  \item \textbf{Batched evaluation}: all evaluation maps are played in lockstep with one batched
        network call per step; it reproduces the earlier sequential results exactly.
  \item \textbf{Tests}: 35 unit tests, e.g.\ GAE vs a reference implementation, a PPO bandit learning
        check, and sync vs subprocess env equivalence.
\end{itemize}

%=====================================================================
\section{Likely Questions and Short Answers}
\begin{description}[style=nextline]
  \item[What is PPO in one sentence?]
  An actor-critic policy-gradient method that does several epochs of SGD on each batch of fresh
  experience, while clipping the probability ratio so the new policy cannot move too far from the
  old one.
  \item[Why clip instead of a KL penalty?]
  It is simpler (first-order, no adaptive penalty coefficient) and works as well in practice. We
  still monitor KL and stop early as a safety net.
  \item[Why is PPO on-policy if it reuses data for 4 epochs?]
  The data comes from the policy just before the update. The ratio $r_t$ corrects for the small
  change during the epochs, and clipping keeps that correction valid. Data is thrown away after
  each update.
  \item[What does the critic do?]
  It estimates $V(s)$, which serves as the baseline for advantages (variance reduction) and for
  bootstrapping in GAE.
  \item[What do $\gamma$ and $\lambda$ control?]
  $\gamma$: how far into the future rewards count. $\lambda$: the bias--variance trade-off of the
  advantage estimate.
  \item[Why entropy regularisation?]
  It prevents premature collapse to a deterministic policy, keeping exploration going and the
  policy stochastic, which matters under partial observability.
  \item[Why does DQN need a target network and replay?]
  Without them the targets move with every update and the samples are correlated, which together
  with bootstrapping makes training diverge.
  \item[What problem does Double DQN solve?]
  Overestimation from taking the max over noisy Q-estimates. One network selects the action and
  the other evaluates it.
  \item[Why is independent learning unstable?]
  Each agent's environment includes a learning opponent, so it is non-stationary. This leads to
  chasing, cycling and forgetting.
  \item[How does self-play fix that?]
  The opponent is drawn from a pool of frozen past versions. This gives each update a fixed target
  and keeps old strategies in the training mix, like fictitious play.
  \item[Why a curriculum?]
  On the full map rewards are sparse. Easy levels give a learning signal, and the skills transfer
  because the observation is size-independent.
  \item[How do you know the agent generalizes and does not memorise?]
  Separate train/val/test seed ranges plus out-of-distribution larger maps. We report the
  train$-$test gap, and test is never used for tuning.
  \item[What is exploitability here?]
  How quickly a fresh PPO agent learns to beat the frozen model, which approximates the
  best-response value.
  \item[Is the game zero-sum?]
  The win/loss reward is ($+1/-1$). Step penalties and shaping make the total slightly non-zero-sum.
  \item[Why were the Hunter's advantages (sprint, spawn near exit) needed?]
  With equal speed and symmetric spawns the Evader won about 99\% of games, so there was no
  competition to learn from. The handicaps balance the game.
  \item[What is reward shaping and is it safe?]
  Extra dense reward for progress (0.01 per tile). Potential-based shaping preserves optimal
  policies; ours is a distance-difference, close to potential-based, and small relative to $\pm1$.
\end{description}

%=====================================================================
\section{Limitations and Future Work}
\begin{itemize}
  \item Few seeds and a short training budget, so results are preliminary.
  \item The memoryless MLP policy is weak under partial observability; an LSTM/GRU policy would help.
  \item The pool samples opponents uniformly; prioritised fictitious self-play (sampling opponents
        we lose to) would be stronger.
  \item Exploitability is approximated by a single fresh PPO learner with a fixed budget.
  \item A fully vectorized (array-based) environment would allow GPU acceleration.
\end{itemize}

\section*{References}
\begin{enumerate}[label={[\arabic*]},leftmargin=*]
  \item J. Schulman et al. \emph{Proximal Policy Optimization Algorithms.} arXiv:1707.06347, 2017.
  \item J. Schulman et al. \emph{High-Dimensional Continuous Control Using Generalized Advantage Estimation.} ICLR, 2016.
  \item J. Schulman et al. \emph{Trust Region Policy Optimization.} ICML, 2015.
  \item V. Mnih et al. \emph{Human-level control through deep reinforcement learning.} Nature, 2015.
  \item H. van Hasselt, A. Guez, D. Silver. \emph{Deep Reinforcement Learning with Double Q-learning.} AAAI, 2016.
  \item Y. Bengio et al. \emph{Curriculum Learning.} ICML, 2009.
  \item J. Heinrich, D. Silver. \emph{Deep Reinforcement Learning from Self-Play in Imperfect-Information Games.} 2016.
  \item O. Vinyals et al. \emph{Grandmaster level in StarCraft II using multi-agent RL.} Nature, 2019.
  \item K. Cobbe et al. \emph{Leveraging Procedural Generation to Benchmark Reinforcement Learning.} 2019.
  \item S. Singh, T. Jaakkola, M. Jordan. \emph{Learning Without State-Estimation in Partially Observable Markovian Decision Processes.} ICML, 1994.
  \item A. Ng, D. Harada, S. Russell. \emph{Policy invariance under reward transformations.} ICML, 1999.
  \item R. Sutton, A. Barto. \emph{Reinforcement Learning: An Introduction.} 2nd ed., 2018.
\end{enumerate}

\end{document}
